%% Preamble — a_linear_algebra_primer_for_financial_engineering
%% Generated in Phase 0 from the scanned source.

\documentclass[11pt,openany]{book}

\usepackage[T1]{fontenc}
\usepackage[utf8]{inputenc}
\usepackage{lmodern}
\usepackage{microtype}

%% -------------------------------------------------------
%% Page geometry
%%
%% The scan measures 6.09in x 9.01in with a 4.69in (338.9pt) text block.
%% That measure is too narrow for this book's own display mathematics: the
%% block-matrix equations of Chapters 4, 6 and 7 run up to 524pt wide and
%% spilled into the right margin. Enlarged to a standard 7x10in textbook
%% trim, which widens the measure to 5.60in (404.5pt, +19%) and deepens it
%% to 8.30in (+12%) — roughly a third more content per page. Base size
%% stays at 11pt so the line still runs about 80 characters.
%%
%% Displays wider than 404.5pt are still broken by hand, following the
%% line breaks the printed book uses; \shrinkwrap below is the last resort.
%% -------------------------------------------------------
\usepackage[
  paperwidth=7.00in,
  paperheight=10.00in,
  left=0.70in,
  right=0.70in,
  top=0.80in,
  bottom=0.90in,
  headheight=14pt,
  headsep=18pt,
  footskip=24pt
]{geometry}

%% -------------------------------------------------------
%% Math
%% -------------------------------------------------------
\usepackage{amsmath}
\usepackage{amssymb}
\usepackage{amsthm}
\usepackage{mathtools}
\usepackage{bm}

%% -------------------------------------------------------
%% Graphics and floats
%% -------------------------------------------------------
\usepackage{graphicx}
\usepackage{float}          % provides the [H] float specifier
\usepackage{adjustbox}
\graphicspath{{./figures/}{../figures/}}

%% Last-resort fit for a single display that no line break can rescue.
%% Scales the body down only if it exceeds the measure; leaves it untouched
%% otherwise, so it never enlarges type. Prefer breaking the display by hand.
%%   \[ \shrinkwrap{ ...matrices... } \]
\newcommand{\shrinkwrap}[1]{%
  \adjustbox{max width=\linewidth}{$\displaystyle #1$}}

%% -------------------------------------------------------
%% Lists, tables, index
%% -------------------------------------------------------
\usepackage[shortlabels]{enumitem}
\usepackage{array}
\usepackage{booktabs}
\usepackage{makeidx}
\makeindex

%% -------------------------------------------------------
%% Links
%% -------------------------------------------------------
\usepackage{xcolor}
\usepackage[
  colorlinks=true,
  linkcolor={blue!60!black},
  citecolor={blue!60!black},
  urlcolor={blue!60!black}
]{hyperref}

%% -------------------------------------------------------
%% Headers / footers — running heads match the source:
%% verso "N.M. SECTION TITLE", recto "CHAPTER N. CHAPTER TITLE"
%% -------------------------------------------------------
\usepackage{fancyhdr}
\pagestyle{fancy}
\fancyhf{}
\fancyhead[LE]{\small\thepage\quad\nouppercase{\leftmark}}
\fancyhead[RO]{\small\nouppercase{\rightmark}\quad\thepage}
\renewcommand{\headrulewidth}{0pt}

%% -------------------------------------------------------
%% Spacing
%% -------------------------------------------------------
\setlength{\parindent}{1.2em}
\setlength{\parskip}{0pt}
\raggedbottom
\emergencystretch=2em

%% -------------------------------------------------------
%% Numbering — equations, figures and tables run per chapter
%% -------------------------------------------------------
\numberwithin{equation}{chapter}
\numberwithin{figure}{chapter}
\numberwithin{table}{chapter}
\setcounter{tocdepth}{2}
\setcounter{secnumdepth}{3}

%% The wider measure lets some long List-of-Tables captions reach the folio
%% column, where book.cls sets them flush against the page number. Reserving
%% more room on the right makes those entries wrap instead of colliding.
\makeatletter
\renewcommand{\@pnumwidth}{2.1em}
\renewcommand{\@tocrmarg}{3.6em}
\makeatother

%% -------------------------------------------------------
%% Theorem environments
%%
%% IMPORTANT — this source uses SEPARATE counters per environment type,
%% each reset per chapter. This is the documented exception to the usual
%% shared-counter rule, and it is verified on the scanned pages:
%%   p.124  Definition 4.7, Definition 4.8, Theorem 4.6  (same page)
%%   p.113  Definition 4.2 immediately followed by Theorem 4.1
%%   p.140  Theorem 5.1 immediately followed by Lemma 5.1
%% Do NOT convert these to a shared counter — it would renumber the book.
%% -------------------------------------------------------
\theoremstyle{plain}
\newtheorem{theorem}{Theorem}[chapter]
\newtheorem{lemma}{Lemma}[chapter]
\newtheorem{corollary}{Corollary}[chapter]
\newtheorem{proposition}{Proposition}[chapter]

\theoremstyle{definition}
\newtheorem{definition}{Definition}[chapter]

%% The source sets "Example:", "Answer:" and "Remark:" as unnumbered runin
%% italic labels, and closes them with the same open square as a proof.
\newenvironment{examplex}%
  {\par\medskip\noindent\textit{Example:}\ }%
  {\hfill$\square$\par\medskip}
\newenvironment{answerx}%
  {\par\medskip\noindent\textit{Answer:}\ }%
  {\par\medskip}
\newenvironment{remarkx}%
  {\par\medskip\noindent\textit{Remark:}\ }%
  {\par\medskip}

%% The proof environment auto-appends \qedsymbol.
%% NEVER write \qed, \blacksquare or \hfill$\blacksquare$ inside a proof.
%% Use \qedhere only when a proof ends with a displayed equation or list.
\renewcommand{\qedsymbol}{$\square$}

%% -------------------------------------------------------
%% Exercises — plain arabic list, restarting at 1 in each chapter's
%% "N.M Exercises" section. Sub-parts are lowercase roman in parentheses.
%% -------------------------------------------------------
\newcommand{\exerciselabel}{\arabic*.}
\newenvironment{exercises}%
  {\begin{enumerate}[leftmargin=*,itemsep=0.9em,label=\exerciselabel,ref=\arabic*]}%
  {\end{enumerate}}
\newenvironment{subparts}%
  {\begin{enumerate}[leftmargin=*,itemsep=0.4em,label=(\roman*),ref=(\roman*)]}%
  {\end{enumerate}}

%% -------------------------------------------------------
%% Metadata macros (title / author / edition only — never a publisher)
%% -------------------------------------------------------
\newcommand{\booktitle}{A Linear Algebra Primer for Financial Engineering}
\newcommand{\booksubtitle}{Covariance Matrices, Eigenvectors, OLS, and more}
\newcommand{\bookauthor}{Dan Stefanica}
\newcommand{\bookedition}{1}

%% -------------------------------------------------------
%% Number fields and script letters
%% -------------------------------------------------------
\newcommand{\R}{\mathbb{R}}
\newcommand{\C}{\mathbb{C}}
\newcommand{\N}{\mathbb{N}}
\newcommand{\Z}{\mathbb{Z}}
\newcommand{\Q}{\mathbb{Q}}
\newcommand{\Prob}{\mathbb{P}}
\newcommand{\Expect}{\mathbb{E}}

\newcommand{\calA}{\mathcal{A}}
\newcommand{\calB}{\mathcal{B}}
\newcommand{\calD}{\mathcal{D}}
\newcommand{\calF}{\mathcal{F}}
\newcommand{\calG}{\mathcal{G}}
\newcommand{\calL}{\mathcal{L}}
\newcommand{\calN}{\mathcal{N}}
\newcommand{\calO}{\mathcal{O}}
\newcommand{\calP}{\mathcal{P}}
\newcommand{\calQ}{\mathcal{Q}}
\newcommand{\calS}{\mathcal{S}}
\newcommand{\calV}{\mathcal{V}}
\newcommand{\calX}{\mathcal{X}}

%% -------------------------------------------------------
%% Probability / statistics operators
%% The source prints var, cov and corr in lowercase upright type and the
%% expectation as an italic capital E; both cases are provided so that
%% either spelling an agent reaches for resolves.
%% -------------------------------------------------------
\newcommand{\E}[1]{E\!\left[#1\right]}
\newcommand{\Var}{\operatorname{var}}
\newcommand{\var}{\operatorname{var}}
\newcommand{\Cov}{\operatorname{cov}}
\newcommand{\cov}{\operatorname{cov}}
\newcommand{\Corr}{\operatorname{corr}}
\newcommand{\corr}{\operatorname{corr}}
\newcommand{\Std}{\operatorname{std}}
\newcommand{\std}{\operatorname{std}}
\DeclareMathOperator*{\plim}{plim}
\newcommand{\dto}{\xrightarrow{d}}
\newcommand{\pto}{\xrightarrow{p}}
\newcommand{\asto}{\xrightarrow{\mathrm{a.s.}}}
\newcommand{\iid}{\overset{\mathrm{iid}}{\sim}}
\newcommand{\Normal}{\mathcal{N}}
\newcommand{\VaR}{\operatorname{VaR}}

%% -------------------------------------------------------
%% Linear algebra operators used by this book
%% -------------------------------------------------------
\DeclareMathOperator{\rank}{rank}
\DeclareMathOperator{\Null}{Null}
\DeclareMathOperator{\Range}{Range}
\DeclareMathOperator{\Span}{span}
\DeclareMathOperator{\diag}{diag}
\DeclareMathOperator{\tr}{tr}
\DeclareMathOperator{\dete}{det}
\DeclareMathOperator{\sgn}{sgn}
\DeclareMathOperator*{\argmax}{arg\,max}
\DeclareMathOperator*{\argmin}{arg\,min}

%% The source writes the transpose as a lowercase superscript t: A^t
\newcommand{\tp}{^{t}}
\newcommand{\transpose}{^{t}}
\newcommand{\inv}{^{-1}}

%% Calculus
\newcommand{\pd}[2]{\frac{\partial #1}{\partial #2}}
\newcommand{\pdd}[2]{\frac{\partial^{2} #1}{\partial #2^{2}}}
\newcommand{\dd}[2]{\frac{d #1}{d #2}}
\newcommand{\grad}{\nabla}

%% Delimiters
\newcommand{\abs}[1]{\left|#1\right|}
\newcommand{\norm}[1]{\left\|#1\right\|}
\newcommand{\inner}[2]{\left(#1,#2\right)}
\newcommand{\set}[1]{\left\{#1\right\}}

%% -------------------------------------------------------
%% Bold vector and matrix macros.
%%
%% The only symbol this book actually sets in bold is the vector of ones
%% (\bfone), e.g. $\bfone^{t} w_{min}$ in Chapter 9. The remaining macros
%% are defined anyway: an undefined \bf<Name> is a fatal error that
%% cascades through an entire chapter, so the whole alphabet plus the
%% Greek letters are provided up front. Never use \vec{} or a raw
%% \mathbf{} in the chapter sources — always a \bf<Name> macro.
%% -------------------------------------------------------
\newcommand{\bfzero}{\mathbf{0}}
\newcommand{\bfone}{\mathbf{1}}

\newcommand{\bfa}{\mathbf{a}}   \newcommand{\bfb}{\mathbf{b}}
\newcommand{\bfc}{\mathbf{c}}   \newcommand{\bfd}{\mathbf{d}}
\newcommand{\bfe}{\mathbf{e}}   \newcommand{\bff}{\mathbf{f}}
\newcommand{\bfg}{\mathbf{g}}   \newcommand{\bfh}{\mathbf{h}}
\newcommand{\bfi}{\mathbf{i}}   \newcommand{\bfj}{\mathbf{j}}
\newcommand{\bfk}{\mathbf{k}}   \newcommand{\bfl}{\mathbf{l}}
\newcommand{\bfm}{\mathbf{m}}   \newcommand{\bfn}{\mathbf{n}}
\newcommand{\bfo}{\mathbf{o}}   \newcommand{\bfp}{\mathbf{p}}
\newcommand{\bfq}{\mathbf{q}}   \newcommand{\bfr}{\mathbf{r}}
\newcommand{\bfs}{\mathbf{s}}   \newcommand{\bft}{\mathbf{t}}
\newcommand{\bfu}{\mathbf{u}}   \newcommand{\bfv}{\mathbf{v}}
\newcommand{\bfw}{\mathbf{w}}   \newcommand{\bfx}{\mathbf{x}}
\newcommand{\bfy}{\mathbf{y}}   \newcommand{\bfz}{\mathbf{z}}

\newcommand{\bfA}{\mathbf{A}}   \newcommand{\bfB}{\mathbf{B}}
\newcommand{\bfC}{\mathbf{C}}   \newcommand{\bfD}{\mathbf{D}}
\newcommand{\bfE}{\mathbf{E}}   \newcommand{\bfF}{\mathbf{F}}
\newcommand{\bfG}{\mathbf{G}}   \newcommand{\bfH}{\mathbf{H}}
\newcommand{\bfI}{\mathbf{I}}   \newcommand{\bfJ}{\mathbf{J}}
\newcommand{\bfK}{\mathbf{K}}   \newcommand{\bfL}{\mathbf{L}}
\newcommand{\bfM}{\mathbf{M}}   \newcommand{\bfN}{\mathbf{N}}
\newcommand{\bfO}{\mathbf{O}}   \newcommand{\bfP}{\mathbf{P}}
\newcommand{\bfQ}{\mathbf{Q}}   \newcommand{\bfR}{\mathbf{R}}
\newcommand{\bfS}{\mathbf{S}}   \newcommand{\bfT}{\mathbf{T}}
\newcommand{\bfU}{\mathbf{U}}   \newcommand{\bfV}{\mathbf{V}}
\newcommand{\bfW}{\mathbf{W}}   \newcommand{\bfX}{\mathbf{X}}
\newcommand{\bfY}{\mathbf{Y}}   \newcommand{\bfZ}{\mathbf{Z}}

\newcommand{\bfalpha}{\boldsymbol{\alpha}}
\newcommand{\bfbeta}{\boldsymbol{\beta}}
\newcommand{\bfgamma}{\boldsymbol{\gamma}}
\newcommand{\bfdelta}{\boldsymbol{\delta}}
\newcommand{\bfepsilon}{\boldsymbol{\epsilon}}
\newcommand{\bfvarepsilon}{\boldsymbol{\varepsilon}}
\newcommand{\bfzeta}{\boldsymbol{\zeta}}
\newcommand{\bfeta}{\boldsymbol{\eta}}
\newcommand{\bftheta}{\boldsymbol{\theta}}
\newcommand{\bfiota}{\boldsymbol{\iota}}
\newcommand{\bfkappa}{\boldsymbol{\kappa}}
\newcommand{\bflambda}{\boldsymbol{\lambda}}
\newcommand{\bfmu}{\boldsymbol{\mu}}
\newcommand{\bfnu}{\boldsymbol{\nu}}
\newcommand{\bfxi}{\boldsymbol{\xi}}
\newcommand{\bfpi}{\boldsymbol{\pi}}
\newcommand{\bfrho}{\boldsymbol{\rho}}
\newcommand{\bfsigma}{\boldsymbol{\sigma}}
\newcommand{\bftau}{\boldsymbol{\tau}}
\newcommand{\bfupsilon}{\boldsymbol{\upsilon}}
\newcommand{\bfphi}{\boldsymbol{\phi}}
\newcommand{\bfvarphi}{\boldsymbol{\varphi}}
\newcommand{\bfchi}{\boldsymbol{\chi}}
\newcommand{\bfpsi}{\boldsymbol{\psi}}
\newcommand{\bfomega}{\boldsymbol{\omega}}

\newcommand{\bfGamma}{\boldsymbol{\Gamma}}
\newcommand{\bfDelta}{\boldsymbol{\Delta}}
\newcommand{\bfTheta}{\boldsymbol{\Theta}}
\newcommand{\bfLambda}{\boldsymbol{\Lambda}}
\newcommand{\bfXi}{\boldsymbol{\Xi}}
\newcommand{\bfPi}{\boldsymbol{\Pi}}
\newcommand{\bfSigma}{\boldsymbol{\Sigma}}
\newcommand{\bfUpsilon}{\boldsymbol{\Upsilon}}
\newcommand{\bfPhi}{\boldsymbol{\Phi}}
\newcommand{\bfPsi}{\boldsymbol{\Psi}}
\newcommand{\bfOmega}{\boldsymbol{\Omega}}
